% Clase 27 --- Doppler con el OBSERVADOR en movimiento (§1.2).
% La fuente esta quieta: los frentes son circulos concentricos EQUIESPACIADOS,
% separados lambda. Lo unico que cambia es cuantos frentes barre el observador
% por unidad de tiempo, porque el va al encuentro de ellos.
\begin{center}
\begin{tikzpicture}[scale=1]

  % fuente en reposo
  \filldraw[figamber] (0,0) circle (2.4pt);
  \node[figlbl, figamber, anchor=north, fill=white, inner sep=1.5pt]
    at (0,-0.16) {fuente en reposo};

  % frentes de onda: circulos equiespaciados
  \foreach \rr in {0.7,1.4,2.1,2.8}{\draw[figaux] (0,0) circle (\rr);}
  \draw[figgray, line width=0.7pt,
        {Latex[length=1.6mm,width=1.2mm]}-{Latex[length=1.6mm,width=1.2mm]}]
    (1.4,0.3) -- (2.1,0.3);
  \node[figlblaux, anchor=south, inner sep=1.5pt, fill=white]
    at (1.75,0.32) {$\lambda$};

  % el observador viene hacia la fuente
  \filldraw[figred] (4.1,0) circle (2.4pt);
  \draw[figvecres] (3.95,0) -- (3.0,0);
  \node[figlbl, figred, anchor=north] at (3.5,-0.12) {$v_0$};
  \node[figlbl, figred, anchor=south] at (4.1,0.2) {observador};

  \node[figlbl, anchor=north, align=center] at (0.6,-3.15)
    {en reposo cuenta $v\,t/\lambda$ frentes
     $\Rightarrow f = v/\lambda$\\[2pt]
     yendo al encuentro barre $(v+v_0)\,t/\lambda$
     $\Rightarrow f' = f\left(1 + v_0/v\right)$};

\end{tikzpicture}\\[0.5em]
{\small\itshape Medir una frecuencia es, literalmente, \textbf{contar frentes de
onda} en un tiempo largo y dividir. Con todo quieto pasan $v\,t/\lambda$ frentes
y sale $f = v/\lambda$. Si el observador se mueve hacia la fuente, los frentes
siguen dibujados exactamente igual ---la fuente no se enteró de nada, así que
$\lambda$ \textbf{no cambia}---, pero él los va a buscar: en el tiempo $t$ el
encuentro cubre $(v + v_0)t$ en vez de $v\,t$, y por eso cuenta más. Sólo
interviene la componente de $v_0$ \textbf{en la dirección de propagación}: quien
se mueve tangencialmente, sobre un mismo círculo, no percibe cambio alguno. Es
importante retener qué es lo que cambia en este caso ---el ritmo del conteo, no
la separación de los frentes---, porque en el caso de la fuente en movimiento
ocurre justo lo contrario, y de esa diferencia sale que las dos fórmulas no
coincidan.}
\end{center}
